\section{Related Literature} \label{literature}

%This paper bridges three strands of literature: (i) responsibility avoidance through delegation, an established line of experimental work in which the intermediary is another human; (ii) algorithm use, which examines when and why people are willing to use algorithmic decision aids, for themselves and others; and (iii) responsibility attributions to decisions involving algorithms, an emerging literature at the intersection of the first two.

%[Comment: Is this first paragraph really necessary? Is it perhaps better placed in a conclusionary paragraph of the whole literature section, perhaps introducing the summary of the contribution of the paper (e.g. in the intro)?]

\subsection*{Responsibility avoidance through delegation}

Delegation has long been studied as an efficiency-enhancing institutional arrangement: principals delegate to agents who have better information, lower opportunity costs, or specialised skills. Alongside these efficiency motives, a separate literature has identified a range of strategic motives for delegation, e.g. as a commitment device that alters the equilibrium (Fershtman and Gneezy, 2001), as a way to motivate the agent (Aghion and Tirole, 1997), or as a means of shifting responsibility for the outcome away from the principal, i.e. principals may delegate not because the agent is more competent or cheaper, but because doing so shifts responsibility for the resulting outcome away from themselves --- both in their own eyes (felt responsibility for the outcome) and in the eyes of others (the blame others attribute to them). The political-science and public-choice literatures have long recognized that delegation can serve as a strategy for shifting blame onto an intermediary~\citep{weaver_politics_1986, fiorina_legislator_1986}. Experimental economics has established the same motive in incentivised settings, primarily using modified dictator games.

\citet{bartling_shifting_2012} provide the canonical experimental evidence. In their design, a dictator may either choose between a fair and an unfair allocation of an endowment or instead delegate this choice to another person whose monetary incentives are aligned with the dictator's. The recipient is adversely affected if the unfair allocation is chosen and may punish the dictator and the intermediary. They find that through delegation punishment is effectively shifted from the dictator to the intermediary, and the prospect of avoiding punishment constitutes a strong motive for the delegation decision itself. \citet{coffman_intermediation_2011} reinterprets the mechanism behind this punishment reduction. He finds that the principal can no longer avoid punishment when she directly interacts with the receiver despite the intermediary's presence, suggesting that what reduces punishment is the absence of direct interaction with the affected party rather than any transfer of blame to a deciding intermediary. This holds even when the intermediary remains completely passive. \citet{oexl_shifting_2013} sharpen this reading by modifying the canonical game so that delegation necessarily eliminates the fair option. The principal nevertheless escapes punishment, with the intermediary actively choosing between two unfair allocations after delegation. Finally, both \citet{bartling_shifting_2012} (with a die roll) and \citet{oexl_shifting_2013} (with an intermediary-triggered random draw) include treatments that substitute a randomization device for the human intermediary. In both, the principal still shifts some punishment to the device, though less than under delegation to a deciding human. Across the two papers, two distinct components of the responsibility-shifting effect emerge --- the absence of direct interaction between principal and recipient drives the reduction in principal punishment, while the act of choosing, rather than any contribution to outcome unfairness, is what transfers blame to the intermediary.

The motive is not confined to settings with third-party punishment.  \citet{hamman_self-interest_2010} extend this picture in a setting without third-party punishment. Still, principals' transfers to recipients fall sharply when they delegate, indicating that responsibility-shifting can operate through felt responsibility or moral wiggle room rather than through the avoidance of punishment. Qualitative self-reported responsibility and recipient-side blame attributions are consistent with this reading---dictators feel less responsible when they delegate, and recipients place less blame on principals than on intermediaries.

%Lastly, \citet{argenton_potters_yang_2023} document the mirror pattern on the credit side: delegators receive less credit for fair outcomes when they delegate than when they choose directly, so that the responsibility-attribution structure operates symmetrically across blame and credit---though magnitudes are not.

Beyond the canonical dictator-game setting, the responsibility-shifting motive --- and the strategic use of delegation it gives rise to --- has been documented across a range of domains. In incentivised ultimatum bargaining, proposers using a hired delegate extract larger surpluses \citep{fershtman_gneezy_2001}. In sequential voting, non-pivotal voters are punished less than pivotal ones for the same group outcome and voters strategically vote against their preferred allocation to avoid being pivotal \citep{bartling_fischbacher_schudy_2015}. For acts of deception, principals pay an agent to lie on their behalf \citep{erat_avoiding_2013}. In policy-making, observers rate delegating legislatures as less blameworthy than legislatures who decide directly, even when the delegate is powerless to change the outcome \citep{hill_does_2015}. \citet{freer_friedman_weidenholzer_2024} isolate the motive in financial decision-making, finding that a substantial fraction of investors delegate even purely luck-based lottery choices, where alternative motives --- decision costs, performance-chasing, or risk tolerance --- cannot rationalize the choice. 

%\citet{steffel_passing_2016} show that people are more willing to delegate decisions made on behalf of others than decisions made for themselves, especially when those decisions might have negative outcomes. Thus, people seem to care more about avoiding blame for bad outcomes than about claiming credit for good ones. [Comment: Perhaps this can be left out here, because we come back to it in the results section?]

In sum, the canonical literature documents that in settings with self-interested choice, delegation reduces principal punishment and her felt responsibility, in turn motivating delegation. Our setting departs from this canonical design along three dimensions: the intermediary is algorithmic rather than human, the outcome falls entirely on a third party rather than partly on the principal herself, and the underlying decision involves genuine predictive uncertainty rather than a known fair/unfair choice. Whether the responsibility-shifting motive retains its force under these conditions is the question we take up.

\subsection*{Delegation to algorithms}
A large and rapidly growing empirical literature studies when and why people are willing to use algorithmic decision aids, beginning with \citet{dietvorst_algorithm_2015}, who introduce the term \textit{algorithm aversion} --- the finding that people prefer their own (or other humans') judgement over algorithmic predictions, even after observing that algorithms outperform human decision-makers. In contrast, \citet{logg_algorithm_2019} document conditions under which subjects show \textit{appreciation} of algorithmic advice over human advice; the picture that emerges is therefore one of substantial heterogeneity in algorithm preferences across tasks, contexts, and elicitation methods. \citet{burton_systematic_2020}, \citet{mahmud_what_2022}, \citet{chugunova_we_2022}, and \citet{qin_ai_2025} provide systematic reviews of this literature and document a large array of moderators for when and why algorithms are used in decision-making. For example, algorithm aversion is stronger when the algorithm cannot be observed to learn from its mistakes \citep{berger_watch_2021}, when the algorithm's reasoning is opaque \citep{litterscheidt_financial_2020,sharan_dont_2020}, and when the algorithm is presented as more complex \citep{stein_dont_2020}. Algorithm aversion is also stronger when the task is perceived as more uncertain \citep{dietvorst_people_2020}, when delegation removes any ability to override the algorithm's choice \citep{dietvorst_overcoming_2018}, and when the algorithm is used in tasks perceived as subjective \citep{castelo_task-dependent_2019,bogert_humans_2021}. Recent additions include \citet{bockstedt_humans_2025}, who show that loss-framing of decision tasks eliminates the preference for human assistance over AI assistance, and \citet{ivanova_stenzel_measuring_2024}, who develop a menu-based design to elicit willingness to cede control to algorithms versus humans.

[Comment:
- the literature reviews and references in general are old-ish...ideally we can include more recent references.
- Maybe include the Dargies, Hakimov, Kübler reference here.
- Also think about where to include the two ivanova stenzel papers]

\subsection*{Responsibility perceptions of decisions involving algorithms}
Within this body of work, a smaller subset focuses on algorithm-use preferences in settings whose consequences fall partly or wholly on someone other than the decision-maker, sometimes called the \textit{moral domain}. The consensus is that aversion to algorithms is amplified once third-party welfare is at stake. \citet{gogoll_rage_2018} document this in an incentivised calculation task with performance uncertainty (both humans and the algorithm can err), \citet{bigman_people_2018} in vignettes across high-stakes moral domains (driving, legal, medical, military); and \citet{jauernig_people_2022} in an incentivised allocation task that identifies \textit{moral discretion} as the human capacity people specifically value.

Once decisions affect a third party, the responsibility-shifting motive established for human delegation becomes relevant for algorithmic delegation as well.

A first strand of evidence, predominantly vignette-based and unincentivised, examines how third parties attribute responsibility and blame to human actors involved with algorithms. In the medical domain, \citet{pezzo_pezzo_physician_2006} document the same pattern for a doctor who follows a (bad) recommendation from an automated decision aid, and \citet{tobia_when_2021} extend the discussion to liability law: physicians who follow standard-care AI recommendations face reduced legal exposure --- and the authors argue that this constitutes an incentive to follow such recommendations specifically because doing so shields physicians from blame. In the workplace setting, \citet{kurtzberg_attribution_2004} find that observers attribute less responsibility to a company when an accident involves technology. [Comment: This is rather old. I would like to include references only after 2020 for the individual domains] Across multiple moral domains (driving, legal, medical, military), \citet{bigman_people_2018} document a robust ``algorithmic outrage asymmetry'': moral violations by algorithms produce less outrage than identical violations by humans. \citet{shank_attributions_2019} consider moral violations produced by human-algorithm pairs and find that humans who monitor algorithms are blamed less than humans acting alone or humans receiving algorithm recommendations, suggesting that responsibility attributions are sensitive to the precise division of authority between human and machine.

Evidence using incentivised laboratory experiments is younger and smaller. \citet{kirchkamp_sharing_2019} compare a modified dictator game in which the dictator either acts alone or makes the decision jointly with a computer. They find no significant difference in the dictator's self-felt responsibility between human-computer and human-human teams, although a tendency towards more selfish allocations under shared decision-making with humans. \citet{maasland_blame_2022} find in a computer-aided HR setting that algorithm aversion \textit{dominates} blame avoidance: the prospect of using an algorithm to delegate an unpleasant HR decision (a dismissal) does not overcome the underlying aversion to algorithms in this domain.

[Comment: It's hard to follow...there is no real red line here. The first refrence is about comparing human-human to human-algorithm, which is yet another dimension. We are not talking about joint decisions, but full delegation. Nevertheless I think it's nice to include it. Does it not tie into Shank attributions 2019?]

A more recent strand of studies has begun to document a pattern that contrasts with the finding of relative algorithm aversion in the moral domain: \citet{chevrier_algorithm_2024} extend the blame chain to include the programmer behind the algorithm: in their design, AI programmers and the delegator both escape punishment for inegalitarian AI-mediated allocations, and the resulting ``moral wiggle room'' is exploited. \citet{kobis_delegation_2025} document that humans request unethical behaviour from machine agents more readily than from human agents, exploiting AI as a plausible-deniability device. \citet{tontrup_strategic_2025} examine delegation of a dictator-game allocation to ChatGPT and find that delegators transfer less to the recipient than non-delegators, indicating that AI delegation is used strategically to extract surplus while shedding moral responsibility, similarly to delegation to other humans in the canonical literature. \citet{hueholt_trusting_2026} use an incentivised trolley-style donation choice between two real charitable allocations and find that subjects delegate the moral decision to an AI significantly more often than to a human counterpart. Across these four different settings, the AI-as-blame-shield motif is shared, and the direction of the effect on delegation is uniformly positive --- in contrast to \citet{maasland_blame_2022}, where algorithm aversion dominates the responsibility-shifting motive in the same broad domain.

The most closely related study is \citet{feier_hiding_2022}, who study delegation to an algorithm in an incentivised laboratory experiment. In their design, to determine a third party's payoff subjects have the option to either rely on their own performance in a Raven-style logic puzzle or, depending on the treatment,  delegate to another human or an algorithm. The subject's own payoff is unaffected by this choice. Third parties may then reward or punish the principal. The headline finding is that decision-makers are held \textit{less} responsible for decisions delegated to an algorithm than for decisions delegated to another human, conditional on outcome. Interestingly, in their experiment, punishment is only reduced significantly when delegating to an algorithm, but not to another human, compared to self-made decisions. Thus, algorithm-delegation, but not human-delegation, shields the principal from blame for bad outcomes.

Our design extends \citet{feier_hiding_2022} along three dimensions. First, where their algorithm is calibrated at the session level so that the delegation decision reflects beliefs about own ability \textit{relative} to that distribution --- a signal value about Player~A's self-perception that may contaminate downstream punishment as a measure of responsibility shifting per se --- we calibrate the algorithm at the individual level, removing the relative-performance signal of delegation by construction. [Comment: Carefule, this applies to both delegating to the algorithm and delegating to another human.]  Second, where the same subject acts as both delegator and evaluator in their design --- potentially introducing self-justification, hypocrisy aversion, and similar motives into the punishment decision --- we separate the delegator and evaluator roles across distinct subject pools. Third, their design does not include a treatment in which punishment is ruled out by construction, so that delegation rates between the two treatments may be different not because of differences in anticipated punishment, but more generally by preferences for algorithms. Their own analysis shows that the propensity to delegate is driven primarily by self-assessed errors in the logic task, not by the agent type. We add this manipulation directly, isolating the role of anticipated punishment in the delegation decision itself. [Comment: Maybe leave out the second point, because we do not talk about it yet? Otherwise include a footnote that shows empirically that the punishment decisions are linked to whether it was delegated or not...on the other hand this would be expected as well.]

Our experiment is, to our knowledge, the first to our knowledge to provide a clean test of blame-shifting in a setting with inherent uncertainty, while controlling relative-performance concerns. It also tests whether, in such a setting the prospect of punishment itself motivates the delegation decision.

[Comment: Include somewhere about Feier:
- The punishment gap is not specifically outcome-conditional — observers reward AI use in good outcomes as well as bad. This is consistent with a decision-rewarding rather than a blame-shielding reading.
- Regarding the double role of subjects: delegation decision is significantly correlated with punishment difference for good outcomes.]


